\section{The arbitrary-background potential quotient}
\label{sec:r5-arbitrary-background-quotient}

Fix a finite extension $K/\mathbf Q_p$ with its original absolute
value. Let $U$ be an ordinary open in a good Hausdorff boundaryless
quasi-smooth $K$-analytic space of pure dimension $n$.
Throughout this section $D=d'd''$, the test forms are all original
compactly supported CLD smooth complementary forms, and positivity
is the original positive cone used in R. Every positive current
under consideration is symmetric. Strong and weak positivity are
not substituted for that cone.

Write $\mathscr{DSH}_{\mathrm{lb}}(U)$ for the actual real functions
which are locally differences of two members of $\mathcal C_b$.
A function in $\mathscr{PH}(U)$ is a finite continuous function
whose two signs are classical psh. Such a function is locally
bounded. No quotient by a measure-null set occurs in either
definition.

\begin{theorem}[Regularity of the bounded classical curvature kernel]
\label{thm:r5-bounded-kernel-regularity}
If $u,v\in\mathcal C_b(U)$ and $D[u]=D[v]$ on every original
test, then the actual difference $u-v$ is continuous and belongs
to $\mathscr{PH}(U)$. More generally,
\[
 \ker\bigl(D:\mathscr{DSH}_{\mathrm{lb}}(U)
                   \longrightarrow\mathcal D'^{1,1}(U)\bigr)
                         =\mathscr{PH}(U).
\]
The equality is an equality of spaces of actual functions.
It does not assert regularity for the kernel on all locally
integrable functions or on all abstract distributions.
\end{theorem}

\noindent\textbf{Status: \texttt{NEW\_PROVED}} in the stated
bounded-classical scope, relative to the declared internal
inputs and the preceding proofs.

\begin{proof}
Let $H=u-v$. Theorem~\ref{thm:r5-all-extension-curve-harmonicity}
shows that after every complete extension $F/K$, the pullback
of $H$ along every allowed smooth analytic curve is an actual
continuous harmonic function. In particular it has the
algebraic-curve property required by
Theorem~\ref{thm:r5-continuity-from-curves}. That theorem applies
to the actual locally bounded $H$ and proves its ordinary
continuity on $U$.

Both $H$ and $-H$ are therefore upper semicontinuous. To check
the full curve quantifier, let $C\to U_F$ be an arbitrary
$F$-analytic curve map, including curves with boundary or
singularities. Its two pullbacks are continuous. In the
definition of subharmonicity on $C$, test further over any
complete extension $E/F$ by a smooth connected strict
$E$-curve mapping to $C_E$. The composite maps to $U_E$,
so the preceding harmonicity theorem makes both signs
subharmonic on that test curve. Smooth curves can be
localized in good Hausdorff boundaryless ordinary opens,
and subharmonicity is local. Definition 3.4 and \S4.1.3
of \cite{GR25} therefore give subharmonicity on the original
$C$. Thus the full classical psh condition holds for both
signs, including its boundary and nonsmooth curve tests.
They belong to $\mathcal C_b(U)$, and
$H\in\mathscr{PH}(U)$.

For a function in $\mathscr{DSH}_{\mathrm{lb}}$ with zero
curvature, apply this argument on its local difference
presentations. Continuity and classical pluriharmonicity
are local and hence glue. Conversely, for
$H\in\mathscr{PH}$ the already established first-order
positive-current theorem applied to $H$ and $-H$ shows that
$D[H]$ and $-D[H]$ are positive. A symmetric current and its
negative can both be positive only if it is zero: on each
finite moment presentation every symmetric smooth test is
dominated in both signs by a positive control. Thus $D[H]=0$.
\end{proof}

\begin{theorem}[Potential independence for every fixed positive closed background]
\label{thm:r5-arbitrary-background-potential-quotient}
Let $0\leq p<n$, and let $S$ be any symmetric original-positive
$(p,p)$-current on $U$, separately $d'$-closed and $d''$-closed.
For any $u,v\in\mathcal C_b(U)$,
\[
 D[u]=D[v]\quad\Longrightarrow\quad D(uS)=D(vS).
 \tag{Q}
\]
All products by $u$ and $v$ are the actual coefficient-Radon
Borel products. The conclusion holds on every original compact
smooth test of bidegree $(n-p-1,n-p-1)$, including all original
moment refinements.

No smoothness, model origin, analytic-cycle representation, or
pre-existing BT representation is required of $S$.
For $p=n$ the corresponding differentiated statement is vacuous.
\end{theorem}

\noindent\textbf{Status: \texttt{NEW\_PROVED}}. The final
positivity step is a direct use of the previously established
arbitrary-fixed-background first-order construction; the new
input is Theorem~\ref{thm:r5-bounded-kernel-regularity}.

\begin{proof}
Put $H=u-v$. Theorem~\ref{thm:r5-bounded-kernel-regularity}
gives $H,-H\in\mathcal C_b(U)$ as actual functions.
Apply the existing first-order BT construction, or explicitly
Proposition~\ref{prop:r5-prescribed-relative-first-order} to
the B-prescribed rows of $H$ and $-H$. For the arbitrary
fixed background $S$ it gives two positive separately
closed currents
\[
                         D(HS),\qquad D((-H)S).
\]
Borel multiplication and distributional differentiation are
linear, so the second is the negative of the first.
The positive-cone argument at the end of the preceding proof
makes this current zero on every original complementary test.
Linearity now gives $D(uS)-D(vS)=D(HS)=0$.

Only the already defined multiplication of a fixed positive
current by a locally bounded Borel function was used. In
particular no pullback of $S$, no regularization of $S$, and
no moving-background limit occurred.
\end{proof}

\begin{corollary}[Differentiated-slot descent]
\label{cor:r5-differentiated-slot-quotient}
Let $S$ be as above, let $0\leq k\leq n-p$, and let
$u_1,\ldots,u_k,v_1,\ldots,v_k\in\mathcal C_b(U)$.
If $D[u_i]=D[v_i]$ for every $i$, then
\[
              T_{S,k}(u_1,\ldots,u_k)
                    =T_{S,k}(v_1,\ldots,v_k)
\]
on all original tests.
\end{corollary}

\noindent\textbf{Status: \texttt{DIRECT\_COROLLARY}} of (Q)
and the frozen fixed-background BT calculus.

\begin{proof}
Replace one differentiated slot at a time. By symmetry of
the established mixed product, place the slot being changed
first. The product of the remaining slots with $S$ is one
fixed positive separately closed background. Apply (Q) to
this background and the two potentials in that slot.
After the finitely many replacements the displayed equality
follows. The case $k=0$ is the identity $S=S$.
\end{proof}

\begin{definition}[Positive currents with bounded classical local potentials]
\label{def:r5-classical-potential-current}
An original $(1,1)$-current $\alpha$ on $U$ has bounded
classical local potentials if every point has an ordinary
neighborhood $O$ and $a\in\mathcal C_b(O)$ such that
$\alpha|_O=D[a]$.
\end{definition}

These currents are positive and separately closed by B.
The definition is a local-potential condition, and it is not
asserted for all positive closed $(1,1)$-currents.

\begin{corollary}[A current-valued mixed operation with one arbitrary background]
\label{cor:r5-local-potential-current-products}
Let $\alpha_1,\ldots,\alpha_k$ have bounded classical local
potentials on $U$, with $k\leq n-p$. There is a well-defined
original positive separately closed current
\[
                        \alpha_1\wedge\cdots\wedge\alpha_k\wedge S
\]
whose expression on any ordinary open carrying local potentials
$\alpha_i=D[a_i]$ is $T_{S,k}(a_1,\ldots,a_k)$.
It is independent of all local potentials and is symmetric
in the $\alpha_i$. It has the restriction, gluing, additivity,
and positive homogeneity properties inherited from the
bounded-potential BT calculus.
\end{corollary}

\noindent\textbf{Status: \texttt{DIRECT\_COROLLARY}}.

\begin{proof}
Choose simultaneous local potentials on an ordinary cover.
On an overlap, both local potentials of each $\alpha_i$
have the same ambient original curvature current.
Corollary~\ref{cor:r5-differentiated-slot-quotient} gives equality
of the two local mixed currents on every original test.
The ordinary sheaf property for currents glues them uniquely.
Their positivity and separate closedness are local.
Symmetry and the linearity assertions follow from those
properties for the chosen local potentials; potential
independence removes the dependence on those choices.
\end{proof}

\begin{corollary}[Pullback of bounded-potential curvature currents]
\label{cor:r5-potential-current-pullback}
Let $\alpha$ have bounded classical local potentials on an
allowed original $K$-space $X$. Let $L/K$ be any complete
extension and let $F:Z\to X_L$ be analytic with $Z$ allowed
over $L$. The local rule
\[
                      F^*\alpha=D[a_L\circ F],
                \qquad \alpha=D[a]\ \hbox{locally},
\]
defines an original positive separately closed $(1,1)$-current
on $Z$, independent of the chosen original-$K$ potentials.
Iterating the rule with the pullbacks of the original potentials
agrees with further scalar extension and analytic composition.
For maps between allowed original-$K$ spaces, this is the usual
intrinsic functoriality on the stated local-potential class.
\end{corollary}

\noindent\textbf{Status: \texttt{DIRECT\_COROLLARY}} of
Theorem~\ref{thm:r5-all-extension-analytic-kernel}.
This is a pullback for the specified local-potential class,
not for an arbitrary positive current.

\begin{proof}
On each inverse-image neighborhood, the actual function
$a_L\circ F$ has a prescribed smooth positive row pulled
back from its original-$K$ row. Its original curvature is
therefore positive and separately closed by the fixed-row
first-order argument. Two choices of $a$ on an original
overlap have equal ambient curvature; the all-extension
analytic kernel theorem gives equal displayed currents on
the whole inverse-image overlap. They glue uniquely.
Composition and extension commute with the actual
compositions of functions and with their original smooth
rows. Applying the same independence statement to the
original potentials proves the asserted compatibility.
When all intermediate spaces are over the original $K$,
the theorem also applies to every alternative bounded
classical potential on each intermediate space, giving
intrinsic functoriality there. Over an arbitrary new field,
the assertion about iteration is restricted to the induced
original-potential presentations; it is not an extension
of B to all new functions over that field.
\end{proof}

\paragraph{Undifferentiated weights and limits.}
The quotient is for differentiated slots. An undifferentiated
weight is still an actual function: replacing it by another
function with the same curvature need not preserve a weighted
current. For example $D[0]=D[1]=0$, but $0S=0$ and $1S=S$
for every nonzero $S$. The seven established fixed-background
BT properties retain their original functional hypotheses.
This section does not turn arbitrary weak convergence of
curvature currents or arbitrary varying backgrounds into
weighted joint convergence.

\paragraph{Original tests and the remaining boundaries of scope.}
The operation above is an equality of ordinary original
CLD currents. It is not promoted to equality on all delta-test
forms. It does not define the wedge of two completely arbitrary
positive closed currents. An arbitrary permitted $S$ may still
charge an ambient pluripolar set. Neither the regularity
theorem nor the quotient proof requires that such concentration
be excluded.

